%% kapa: exam=1 round=34 year=2023 period=2 form=A subject=accounting subject_ko=회계학 seq=32 number=72
%% source: Q-Net 감정평가사 기출문제 (cst00309) artlSeq=5211789 file=raw/1008/5211789/2023년도 제34회 감정평가사 1차시험 2교시.pdf
%% conversion: kordoc 4.24.51 -> tools/gen1.py (auto); review=proofread-sonnet; flags=none
\begin{question}{72}
(주)감평은 제품 X, Y, Z를 생산·판매하고 있으며, 각 제품 관련 자료는 다음과 같다.
\par{\small\centering\begin{tabular}{|c|c|c|c|}\hline
구 분 & 제품 X & 제품 Y & 제품 Z \\ \hline
매출배합비율(매출수량기준) & 20\% & 60\% & 20\% \\ \hline
단위당 공헌이익 & ￦12 & ￦15 & ￦8 \\ \hline
손익분기점 매출수량 & ? & 7,800단위 & ? \\ \hline
\end{tabular}\par}
(주)감평은 제품 Z의 생산중단을 고려하고 있다. 제품 Z의 생산을 중단하는 경우에 고정비 중 ￦4,000을 회피할 수 있으며, 제품 X와 Y의 매출배합비율(매출수량기준)은 60\%와 40\%로 예상된다. (주)감평이 제품 Z의 생산을 중단할 경우, 목표이익 ￦33,000을 달성하기 위한 제품 X의 매출수량은?
\begin{choices}
  \item 6,900단위
  \item 7,800단위
  \item 8,400단위
  \item 8,700단위
  \item 9,000단위
\end{choices}
\end{question}
\answer{5}
